\section{Results Plan}

This section reports the completed four-dataset RQ1 matched comparison and the
validated RQ1b held-out evaluation, while retaining protocol scaffolds for
RQ2--RQ5. Claims are limited to the immutable artifacts attached to each
subsection.

\subsection{RQ1: effectiveness}
\textbf{Intended claim.} Workload-optimized native statistics can reduce
cardinality-estimation error on some benchmark workloads relative to stock and
simple baselines.

\textbf{Required evidence.} A per-query result table bound to immutable
manifests, plus a distribution plot of q-error and a table of mean, quantiles,
maximum, and query outcome counts.

\begin{figure}[t]
\centering
\fbox{\parbox{0.86\columnwidth}{\centering RQ1 q-error distributions and baseline comparison\par\vspace{0.3em}\small TODO: render from confirmatory derived artifacts}}
\caption{Planned RQ1 figure; no values are supplied in this framework.}
\end{figure}
% The generated cross-dataset table below is derived from
% experiments/rq1-cross-dataset-summary-v2.json (digest
% 84ce7a91fc94ad137f1b8dfd901429d41e3ab89a84c1cece30901233f42da90c).
% This RQ1 comparison is in-workload (the workload used by the advisor
% objective), not a held-out query generalization evaluation. RQ4 heuristic
% baselines are separate and remain planned.

% Metrics are unchanged from the immutable v1 table source; current paper
% provenance is experiments/rq1-cross-dataset-summary-v2.json.
% Summary semantic digest:
% 84ce7a91fc94ad137f1b8dfd901429d41e3ab89a84c1cece30901233f42da90c
\begin{table}[t]
\centering
\scriptsize
\caption{RQ1 three-arm matched comparison on the advisor's in-workload queries. Values are q-error; lower is better. The advisor arm is evaluated after stock PostgreSQL full-data deployment.}
\label{tab:rq1-cross-dataset-summary}
\resizebox{\columnwidth}{!}{%
\begin{tabular}{lrrrrrr}
\toprule
Dataset & Default & Target 10000 & Advisor & Adv. p95 & Adv. p99 & Adv. max \\
\midrule
Census13 & 5.232 & 5.247 & 2.232 & 6.000 & 15.505 & 107.000 \\
Forest10 & 7.144 & 7.150 & 6.669 & 14.603 & 64.000 & 9754.000 \\
Power7 & 127.503 & 130.815 & 85.208 & 13.681 & 161.000 & 157380.000 \\
DMV11 & 276.868 & 282.164 & 41.401 & 28.729 & 996.956 & 27877.133 \\
\bottomrule
\end{tabular}
}%
\end{table}


The four-dataset summary is complete for the declared three-arm, in-workload
comparison under the system freeze recorded by those artifacts. It records
paired classifications, tail-error concentration, and the largest per-query
regressions in addition to the headline aggregates.
The v2 summary also separates observations-file provenance from the
snapshot-bound GroundTruthSet provenance; v1 remains retained as immutable
historical output.
These results do not claim that every query improves, do not complete the RQ4
heuristic study, and do not replace the PostgreSQL estimator with a learned
model.
% TODO(PAPER-LITERATURE): document exact comparability limits for learned-CE context.

\subsection{RQ1b: held-out workload generalization}
Unlike RQ1a, which evaluates a recommendation on the workload used for
selection, RQ1b asks whether native extended statistics selected from a validation
workload retain cardinality-estimation utility on a held-out workload
generated under the same audited AreCEL workload-generator contract. This is a
same-distribution held-out evaluation, not a workload-drift or
distribution-shift experiment. For each dataset, the advisor used only
$W_{\mathrm{valid}}$ and authoritative validation truth to construct and seal
$S_{\mathrm{valid}}$. The sealed recommendation was then deployed to fresh
stock PostgreSQL, and each of the 10,000 test queries was evaluated once with
physical \texttt{EXPLAIN (FORMAT JSON)}. The comparison reuses the immutable
PG16-default and PG16-target10000 baseline arms; it does not rerun either
baseline. The design used 10,000 sampled rows with seed 42, statistics target
100, screening width $K_s=8$, at most $B=8$ selected definitions, and a
300-second search limit.

\begin{table*}[t]
\centering
\small
\caption{RQ1b held-out workload generalization. Statistics are selected using
the validation workload $W_{\mathrm{valid}}$ and evaluated on the sealed
recommendation over the 10,000-query test workload $W_{\mathrm{test}}$ after
deployment to stock PostgreSQL. The three q-error columns report mean
q-error; lower is better. $k$ is the number of definitions in $S_{\mathrm{valid}}$,
and $\Delta$ is the relative mean reduction versus PG16-default.}
\label{tab:rq1b-held-out-generalization}
\begin{tabular}{lrrrrrr}
\toprule
Dataset & $k$ & PG16-default & PG16-target10000 & $S_{\mathrm{valid}}$ & $\Delta$ default & Strict-unseen $n$ \\
\midrule
Census13 & 7 & 5.23 & 5.25 & 2.24 & 57.3\% & 9,386 \\
Forest10 & 7 & 7.14 & 7.15 & 6.68 & 6.5\% & 10,000 \\
Power7 & 7 & 127.50 & 130.81 & 64.36 & 49.5\% & 10,000 \\
DMV11 & 5 & 276.87 & 282.16 & 40.31 & 85.4\% & 8,138 \\
\bottomrule
\end{tabular}
\end{table*}


Table~\ref{tab:rq1b-held-out-generalization} shows that $S_{\mathrm{valid}}$
has lower mean held-out q-error than both historical baseline arms on all four
datasets. The reduction versus PG16-default is 57.3\% for Census13, 6.5\%
for Forest10, 49.5\% for Power7, and 85.4\% for DMV11. Thus the direction is
consistent, but the magnitude is strongly dataset-dependent; the DMV11 value
is not a universal improvement rate. All four searches terminated at the
recorded local optimum, which does not establish global optimality.

The mean reductions are not uniform per-query improvements. Relative to
PG16-default, the paired counts (improved/tied/worsened) are 4,936/2,909/2,155
for Census13, 4,322/2,811/2,867 for Forest10, 5,719/575/3,706 for Power7, and
6,518/1,027/2,455 for DMV11. Consequently, 21.6\%--37.1\% of test instances
worsened despite the lower mean. The p95 and p99 q-errors decrease on all four
datasets, consistent descriptively with fewer severe errors. Maximum q-error
also decreases for Census13, Forest10, and DMV11, but increases slightly for
Power7 (153,148 to 155,314); these tails do not support a causal explanation.

The preregistered strict-unseen analysis filters the same physical test
observations and is secondary, not a second experiment. It removes 614
Census13 instances and 1,862 DMV11 instances whose exact source-query hashes
were seen in validation, leaving 9,386 and 8,138 instances respectively. The
mean reductions versus PG16-default remain positive: 57.95\% for Census13,
84.51\% for DMV11, 6.55\% for Forest10, and 49.52\% for Power7. Forest10 and
Power7 have strict-unseen populations equal to the full test population, so
they provide no within-dataset contrast between those populations. Strict
unseen means exact source-query non-overlap, not a different query
distribution.

The identifier-level intersection between $S_{\mathrm{valid}}$ and the
historical in-workload $S_{\mathrm{test}}$ is zero for all four datasets. The
published evidence does not establish that candidate IDs from independently
constructed snapshots denote comparable physical definitions, so this is not
evidence of disjoint PostgreSQL statistics. In particular, $S_{\mathrm{test}}$
is neither an oracle nor an optimality bound. Census13 also retains its
audited historical baseline caveat: the RQ1a baseline was truth-rebound from
historical evidence, its PG16-default arm inherited target $-1$, and its
PG16-target10000 arm used an explicit target of 10,000; RQ1b selection used
target 100.

\subsection{RQ2: sample-to-full-data transfer}
\textbf{Intended claim.} A design selected on a fixed sample can retain benefit
after native payloads are independently recomputed on full data, subject to
measured transfer error (RQ2a). Separately, the sample-sandbox utility may or
may not predict full-data utility (RQ2b); this is a transfer relationship, not
a mechanism-fidelity claim.

\textbf{Required evidence.} Paired P0/P1/P2 per-query results, ordinary-stat
fingerprints, deployment verification, and the source Recommendation digest.
The analysis will distinguish the fixed-sample objective $J_S(M)$ from the
full-data objective $J_D(M;A)$ for a particular native \texttt{ANALYZE}
realization $A$; native realization variability is not reported as RQ3
mechanism fidelity.
% TODO(PAPER-RQ2): insert P0/P1/P2 transfer table and paired query analysis.
% TODO(PAPER-RQ2B): insert sample-utility/full-data-utility relationship figure.

\subsection{RQ3: fidelity}
\textbf{Intended claim.} Catalogless hypothetical evaluation reproduces
physical PostgreSQL statistics behavior to the extent demonstrated by a
same-patched-binary comparison controlling data, equivalent native payload
semantics, ordinary statistics, and planner/session settings. The primary
evidence is direct root \texttt{Plan Rows} agreement and classified mismatches;
objective difference is derived. A patched-physical versus stock-physical
comparison is a secondary overlay-inactive sanity check, not the primary
fidelity result. The Priority-0 synthetic artifact covers MCV-only,
dependencies-only, and combined MCV-plus-dependencies configurations.
The formal primary artifact reports 9 paired queries, 9 exact matches, and zero
mismatches across the three configurations; the independent secondary artifact
also reports 9 of 9 exact stock-versus-patched physical pairs with the
hypothetical overlay inactive. These results support only the declared
same-binary synthetic mechanism claim and its registered controls.
% TODO(PAPER-RQ3): render tab-rq3-mismatch-classification from the immutable
% primary and secondary artifacts; retain the table as the primary presentation
% if the nine paired observations are clearer than a figure.
% The sample-to-full-data utility relationship is reported under RQ2.

\subsection{RQ4: ablation and advisor necessity}
\textbf{Intended claim.} Planner-in-the-loop search contributes measurable
value beyond low-cost selection rules on the tested candidate universes; the
result need not imply global optimality. RQ4a is the primary deterministic
patched-sandbox design comparison. RQ4b is a secondary stock physical
consequence study and is explicitly subject to native \texttt{ANALYZE}
sampling variability.
\textbf{Required evidence.} RQ4a requires exact deterministic replay of
membership, planner estimates, objective, and configuration traces, plus the
fixed-$k$ method comparison and tiny-universe exhaustive diagnostic. RQ4b
requires a shared stock realization ID, one recorded union \texttt{ANALYZE},
post-drop payload equality without another \texttt{ANALYZE}, and per-method
physical EXPLAIN/q-error evidence. Neither sandbox objective nor one native
realization is substituted for the other layer.
\textbf{Partial dataset-scoped evidence.} The Forest10 fixed-$k$ canary has
passed the deterministic patched-sandbox replay and the shared stock physical
realization gates for nine configurations (five pre-registered random seeds,
workload-frequency, native-payload-size, singleton-utility, and greedy
ADD). The immutable parent artifact is
\texttt{experiments/arecel-forest10/rq4-fixed-k/rq4-ablation-v1.json} (digest
\texttt{736ba74f...}); the global RQ4 status remains incomplete because
Census13, Power7, and DMV11 are still planned. This result is protocol
evidence, not a claim that greedy search is globally optimal or superior.
% TODO(PAPER-RQ4): insert the complete cross-dataset heuristic and
% tiny-universe exhaustive-search table after the remaining datasets run.

\subsection{RQ5: practicality}
\textbf{Intended claim.} The operational cost of offline capture, native
statistics construction, sandbox search, and deployment is measurable and
should be considered alongside accuracy.
% TODO(PAPER-RQ5): insert cost breakdown, storage, refresh, and planning-overhead figures.

\subsection{Pilot evidence ledger}

The repository contains historical pilot artifacts for all four datasets and
validated RQ1 canonical artifacts for Census13, Forest10, Power7, and DMV11
under \texttt{experiments/}. The historical outputs remain preliminary; the
canonical artifacts retain their embedded system-freeze-v1 provenance and are
not retroactively reassigned to system-freeze-v2. The DMV11 full-data
transfer artifact remains separate RQ2 evidence and is not part of the RQ1
completion claim.
% TODO(PAPER-PILOT): link each figure/table to a compact derived artifact digest.
